\section{引言}

\subsection{研究背景}

\subsubsection{人口老龄化与社区养老服务压力}

截至2025年底，我国60周岁及以上老年人口已达3.23亿，占人口总数的23.0\%；65周岁以上人口达2.24亿，占人口总数的15.9\%，全国人均预期寿命约79岁【1】【3】。由此可见，人口老龄化问题是我国当前与今后长时期内较为突出的社会结构性挑战之一。2026年《政府工作报告》中指出，要加强社会保障和服务，深入实施积极应对人口老龄化国家战略，扩大普惠养老服务供给【2】。随着家庭规模缩小，传统家庭赡养功能弱化，养老需求逐渐从基本生活照料扩展至就医、康复护理、助餐等符合型领域。
 
针对养老服务模式，习近平总书记强调：“要构建居家为基础、社区为依托、机构为补充、医养相结合的养老服务体系，更好满足老年人养老服务需求。”【4】社区养老更贴近日常生活场景，也更依赖于社区内部及周边生活圈医疗、交通、休闲、生活服务等设施支撑，这对社区养老服务的供给能力提出了更高要求。因此，如何在有限公共资源约束下最大化社区养老服务的空间匹配度和服务效率，成为基层公共服务优化的重要问题。

\subsubsection{福州市社区养老服务的现实基础}

福州市60周岁以上老年人口已达155万，老龄化率达22\%，已经迈入中度老龄化社会【5】，具有开展社区养老服务空间适配研究的典型性。一方面，福州市较早开展社区养老服务改革的探索，持续推进居家社区养老、普惠养老和智慧养老建设，具备一定的政策和实践基础。

另一方面，福州城市空间结构较为多元，主城区内部既有老城区（鼓楼区、台江区）人口密度高、老旧社区多、老年人口相对集聚的现实情况，也有新区（晋安区、马尾区、仓山区）居住空间扩张较快、服务设施持续完善的问题。不同区域在老年人口分布、医疗照护资源、生活服务配套和公共交通条件等方面存在差异，使社区养老服务面临需求集中、供给不均和区域适配不足等现实压力。因此，福州适合作为社区养老服务空间配置与生活权适配的研究对象


\subsubsection{从“设施覆盖”到“有效可达”的研究转向}

在现有社区养老服务研究与实践中，通常以养老设施数量、覆盖范围或服务半径作为衡量供给水平的主要依据。上述指标能够较为直观地反映养老服务资源的总体配置情况，但仍难以充分说明老年人是否能够在日常生活中便捷、持续地获得养老服务【8】。

养老服务嵌入老年人的社区生活圈中，养老服务设施存在并不意味着可达性高；设施覆盖率高也并不等同于供需匹配。【9】单纯依靠设施数量等指标来规划或评价养老服务，极容易忽视老年人在真实日常生活场景中获取养老服务的障碍。

基于以上判断，本文不仅关注福州市主城区社区养老服务设施的空间分布，更进一步评估养老服务与老年人口需求、医疗健康资源、日常生活配套、公共交通条件以及代际共享空间之间的匹配关系，识别社区生活圈中可能存在的养老服务短板与空间错配。本文希望能更真实地刻画老年人综合性社区养老服务获取水平，并基于此为主城区社区养老服务优化与新、旧城区生活圈的发展和更新提出更具针对性的依据

\subsection{研究问题}

\subsubsection{福州市老年服务需求如何空间分布}

回答“老人在哪里”。

\subsubsection{社区养老服务是否与老年需求匹配}

回答“服务是否供给到真正需要的地方”。

\subsubsection{代际共享服务能否缓解养老服务压力}

说明社区食堂、公园、图书馆、社区活动中心、体育空间、便民服务中心等设施虽然不是传统养老设施，但可以为老年群体提供助餐、休闲、社交和生活支持。你们前期材料已经把社区服务划分为银发服务、青年活力和代际共享服务三类，这里可以作为指标体系的来源。

\subsubsection{有限资源下应优先补什么、补在哪里}

回答“补位优化”的问题。

\subsection{研究思路与技术路线}

\subsubsection{总体研究思路}

可以写成：

\begin{quote}
本文按照“需求识别—服务可达—错配诊断—成因解释—智能补位”的逻辑展开。
\end{quote}

\subsubsection{技术路线设计}

建议放一张技术路线图。

图中包括：

数据采集 → 网格化处理 → 老年需求估计 → 代际共享修正 2SFCA → 供需错配识别 → GraphSAGE 风险识别 → LightGBM-SHAP 成因解释 → MCLP 服务补位优化。

\subsubsection{研究创新}

写三点即可。

第一，构建代际共享修正的养老服务可达性测度方法。
第二，引入图神经网络识别空间关联下的高错配风险网格。
第三，将错配识别与 MCLP 补位优化结合，形成可执行的设施布局建议。
